\documentclass[10pt,a4paper]{amsart}
\usepackage[margin=19mm]{geometry}
\usepackage{amsmath,amssymb,mathtools}
\usepackage[scr=rsfs,cal=euler]{mathalfa}
\usepackage{xfrac}
\usepackage[unicode,hidelinks,bookmarks=false]{hyperref}
\newcommand{\ie}{i.e.\ }
\newcommand{\cf}{cf.\ }
\newcommand{\resp}{respectively\ }
\newcommand{\Subsection}{\subsection}
\providecommand{\nobreakdash}{\nobreak\hbox{-}}
\setlength{\emergencystretch}{2em}
\multlinegap0pt
\usepackage{bm}
\usepackage{bbm}
\usepackage{dsfont}
\usepackage{xspace}
\usepackage{cancel}
\usepackage{mathtools}
\usepackage{amsthm}
\makeatletter
\@ifundefined{theorem}{\newtheorem{theorem}{Theorem}}{}
\@ifundefined{corollary}{\newtheorem{corollary}[theorem]{Corollary}}{}
\@ifundefined{definition}{\newtheorem{definition}[theorem]{Definition}}{}
\makeatother
\newcommand{\Cat}{\textrm{Cat}}
\newcommand{\tgt}{\textrm{Tgt}}
\title[Reduction of Probabilistic Chemical Reaction Networks]{Reduction of Probabilistic Chemical Reaction Networks}
\author{Mauricio Montes, Gregoire Sergeant-Perthuis}
\date{Source: arXiv:2606.27737v1 (CC BY 4.0)}
\begin{document}
\maketitle

\noindent\emph{Source and licence.} Reduction of Probabilistic Chemical Reaction Networks. Version arXiv:2606.27737v1 (CC BY 4.0), distributed under the Creative Commons Attribution 4.0 International licence, as recorded in the arXiv licence metadata for this exact version and shown on the abstract page https://arxiv.org/abs/2606.27737v1. Full rights notice: licenses/arxiv-2606.27737-CC-BY-4.0.txt.\par\smallskip

\noindent\textit{Editorial scope.} Counted unit: the Corollary whose source label is \texttt{cor:roundtrip} in the article (source0.tex lines 1776--1813, source label \texttt{cor:roundtrip}), delivered together with the article's own complete original proof of it (source0.tex lines 1815--1882, one \texttt{proof} environment of 68 lines). No mathematics is written, repaired or re-derived: statement and proof are the article's own text line for line. Required context retained uncounted: thm:recon (lines 596--614); def:variable-equivalence (lines 619--637); thm:BP-implementation (lines 672--706). Every definition, display and label of those lines is retained unchanged. Omitted from this package and not redistributed: 1--595, 615--618, 638--671, 707--1775, 1814--1814, 1883--3519. Nothing inside a retained excerpt is omitted or altered: every retained body is delivered exactly as the article's lines are written, so no omission receipt of any kind accompanies this package. Citations: the retained excerpts carry no citation command at all, so no bibliography entry of the article is transcribed and no reference is redistributed. Reference adaptation: none was needed and none was made; every reference inside the delivered unit resolves inside it, so no unresolved reference or citation is produced. Printed slips kept literal: no printed word, symbol, hypothesis or number of the counted statement or of its proof was altered. Layout and class: the article's own class is \texttt{article} with packages including microtype, tikz-cd, graphicx, subfigure, booktabs, xcolor, hyperref, icml2026, amsmath, amssymb, mathtools, amsthm, comment, cleveref; the wrapper is the lane's pinned \texttt{amsart} editorial wrapper at 10pt on A4, which restates the theorem-like environments the retained text needs and adds the article's own macro definitions that the retained text needs (\texttt{\textbackslash Cat}, \texttt{\textbackslash tgt}), with extra wrapper packages (bm, bbm, dsfont, xspace, cancel, mathtools). The delivered numbering is the unit's own. Assets: no retained span contains a figure environment, a graphics-inclusion command, a PDF-page inclusion, a table or a TikZ picture; no image file is copied into this package, the package declares no asset dependency and no image file is redistributed with it. Licence: version arXiv:2606.27737v1 of this article is distributed under the Creative Commons Attribution 4.0 International licence, as recorded in the arXiv licence metadata for this exact version and shown on the abstract page https://arxiv.org/abs/2606.27737v1. A full rights notice is delivered at licenses/arxiv-2606.27737-CC-BY-4.0.txt. Characters: every retained excerpt is pure ASCII, so the delivered engine, XeLaTeX with its default Latin Modern text font, reports no missing character.
\begin{theorem}[CRN $\to$ Factor Graph Reconstruction]\label{thm:recon}
Let $\Gamma = (X, R, \kappa)$ be a mass-action CRN satisfying \textup{(W1)--(W6)}.
Then one can reconstruct a factor graph $\mathrm{FG}(\Gamma) = (V,E, H)$ as follows:
\begin{itemize}
    \item The factor nodes are the factor classes $J$ from \textup{(W6)}.
    \item The variable nodes are the variable classes $[P]$ under $\sim_v$ 
          \textup{(Definition~\ref{def:variable-equivalence})}.
    \item A variable class $[P]$ is incident to a factor class $J$ if and only 
          if $[P] \cap V_J \neq \varnothing$.
    \item The state space of variable class $[P]$ is $E_{[P]} := E_P$ for any 
          representative $P \in [P]$.
    \item The factor table $\psi_J : \prod_{Q \in V_J} E_Q \to (0,\infty)$ from 
          \textup{(W6)} serves as the potential for factor node $J$.
\end{itemize}
Moreover, the incidence relation is well-defined: distinct elements of $V_J$ 
belong to distinct variable classes, establishing a bijection between $V_J$ and 
the variable classes incident to $J$. Each factor table $\psi_J$ is determined 
by $\Gamma$ up to multiplication by a positive constant.
\end{theorem}

\begin{definition}[Variable equivalence on product bundles]\label{def:variable-equivalence}
Assume (W1)--(W6). Let $\mathcal{P}$ be the set of product bundles and 
$\mathcal{S}$ be the set of sum bundles.

Define a binary relation $\sim$ on $\mathcal{P}$ by: for $P, P' \in \mathcal{P}$,
\newline $P \sim P'$ if and only if:
\begin{align*}
    \exists\, B \in \mathcal{S} \text{ such that: } &
    \Bigl[\tgt(B) = P \text{ and } B \in \Cat(P')\Bigr]
     \\ \;\text{ or }\; &
    \Bigl[\tgt(B) = P' \text{ and } B \in \Cat(P)\Bigr].
\end{align*}

Let $\sim_v$ be the equivalence relation on $\mathcal{P}$ generated by $\sim$
(the reflexive, symmetric and transitive closure).

The equivalence classes of $\mathcal{P}$ under $\sim_v$ are called 
\emph{variable classes}.
\end{definition}

\begin{theorem}[BP Implementation]\label{thm:BP-implementation}
Let $\Gamma = (X, R, \kappa)$ be a mass-action CRN satisfying \emph{(W1)--(W6)} and \emph{(R1)}, 
and let $\mathrm{FG}(\Gamma)$ be the factor graph with factor tables $\{\psi_J\}$ reconstructed 
by Theorem~\ref{thm:recon}. Then the positive steady states of $\Gamma$ correspond 
bijectively to BP fixed points on $\mathrm{FG}(\Gamma)$.

More precisely, at any positive steady state. For each factor class $J$ and 
target $Q \in V_J$, write $A_Q = \prod_{Q' \in V_J \setminus \{Q\}} E_{Q'}$ 
for the set of state assignments to the remaining variables, and write 
$\psi_J(k,\, a) := \psi_J\bigl(x_Q = k,\, x_{Q'} = a(Q') \text{ for all } 
Q' \in V_J \setminus \{Q\}\bigr)$ for the corresponding factor-table entry.

\begin{enumerate}
\item For each sum bundle $B = B_{J \to Q}$ associated to factor class $J$ and 
target product bundle $Q \in V_J$, the steady-state concentrations satisfy
\begin{align}\label{eq:sum-steady}
\frac{\rho_B}{[B_0]}[B_k] 
= \kappa_{J \to Q} 
\sum_{a \in A_Q} 
\psi_J(k,\, a)
\prod_{Q' \in V_J\setminus\{Q\}} [Q'_{a(Q')}]
\end{align}
for all $k \in E_Q$.

\item For each product bundle $P$, the steady-state concentrations satisfy
\begin{equation}\label{eq:prod-steady}
\frac{\rho_P}{\kappa_P [P_0]}[P_k] = \prod_{B \in \Cat(P)} [B_k]
\end{equation}
for all $k \in E_P$.

\item These equations are precisely the sum-product BP fixed-point equations on 
$\mathrm{FG}(\Gamma)$, with the factors $\rho_B/[B_0]$ and $\rho_P/(\kappa_P[P_0])$ 
serving as the per-message normalization constants.
\end{enumerate}
\end{theorem}

\begin{corollary}[BP $\leftrightarrow$ CRN]\label{cor:roundtrip}
Let $\Gamma=(X,R,\kappa)$ be a mass--action CRN satisfying
\textup{(W1)--(W6)} and \textup{(R1)}, and let $\mathrm{FG}(\Gamma)$ be the
factor graph with factor tables $\{\psi_J\}$ reconstructed by
Theorem~\textup{\ref{thm:recon}}. Let
\[
  \Gamma^{\mathrm N} \;:=\; \mathcal N\bigl(\mathrm{FG}(\Gamma)\bigr)
\]
be the BP--CRN obtained by applying the Napp--Adams compilation to
$\mathrm{FG}(\Gamma)$, with any fixed choice of recycling and production
constants.

Then there exist positive scalars $\{s_B>0\}$, one for each bundle $B$, and
a choice of rate constants for $\Gamma^{\mathrm N}$ such that the following
hold.
\begin{enumerate}
  \item The species of $\Gamma$ and $\Gamma^{\mathrm N}$ can be put into
        bijection bundlewise, preserving zero versus coordinate species,
        the distinction between sum and product bundles, and catalytic
        neighbor sets, and under this bijection the change of variables
        \[
          [X_{B,k}]^{\mathrm N} \;=\; s_B\,[X_{B,k}]
        \]
        conjugates the mass--action vector field of $\Gamma$ to that of
        $\Gamma^{\mathrm N}$.

  \item In particular, positive steady states of $\Gamma$ correspond
        bijectively, via the same bundlewise rescaling, to positive steady
        states of $\Gamma^{\mathrm N}$.

  \item From the main result of Napp--Adams applied to
        $\Gamma^{\mathrm N}$, any positive steady state of $\Gamma$
        therefore determines a collection of BP messages (defined up to the
        usual per--message positive rescaling) on $\mathrm{FG}(\Gamma)$
        that satisfy the sum--product fixed--point equations for the
        recovered factor tables $\{\psi_J\}$.
\end{enumerate}
\end{corollary}

\begin{proof}[Proof of ~\ref{cor:roundtrip}]
Let $\mathrm{FG}(\Gamma)$ be reconstructed from $\Gamma$ under (W1)--(W6), with factor 
classes $(J, V_J)$ and factor tables $\psi_J$ from (W6). Consider the Napp compilation 
$\Gamma^N := \mathcal{N}(\mathrm{FG}(\Gamma))$, in which each directed incidence 
$(J, Q)$ with $Q \in V_J$ produces a sum bundle $\mathsf{S}^{J \to Q}$ and each 
directed incidence $(Q, J)$ produces a product bundle $\mathsf{P}^{Q \to J}$.

By the construction of $\mathrm{FG}(\Gamma)$ in Theorem~\ref{thm:recon}, 
the bundles of $\Gamma$ decompose into:
\begin{enumerate}
\item Sum bundles $B_{J \to Q}$ indexed by pairs $(J, Q)$ with $Q \in V_J$;
\item Product bundles $P$ indexed by directed incidences $(Q, J)$.
\end{enumerate}

The Napp compilation produces exactly one bundle of each type for each such directed 
incidence. Hence there is a canonical bundlewise bijection between the species sets 
of $\Gamma$ and $\Gamma^N$, preserving bundle type (sum vs.\ product), zero vs.\ 
coordinate species, and state labels.

Moreover, by (W4) and (W5), the reaction monomials of $\Gamma$ match the Napp monomials:
\begin{itemize}
\item Sum-bundle production terms are multilinear monomials with one coordinate 
from each catalytic neighbor;
\item Product-bundle production terms are diagonal in the state label $k$.
\end{itemize}

Fix a factor class $(J, V_J)$ and $Q \in V_J$. By (W6.3), the production coefficients 
of the corresponding sum bundle $B_{J \to Q}$ satisfy:
\[
C_{B_{J \to Q}}(k; a) = \kappa_{J \to Q} \cdot \psi_J(x_Q = k, x_{Q'} = a(Q') \text{ for } Q' \neq Q).
\]

In the Napp compilation, the analogous producing reactions for $\mathsf{S}^{J \to Q}_k$ 
have rate constants proportional to the same clamped table entries of $\psi_J$, with 
a user-chosen proportionality constant. Choose these Napp proportionality constants 
so that, after a bundlewise rescaling
\[
[X_{B,k}]^N = s_B [X_{B,k}],
\]
the coefficients of every monomial in the sum-bundle production polynomials agree 
term-by-term between $\Gamma$ and $\Gamma^N$.

Similarly, by (R1) we may choose Napp recycling and product-production constants to 
match the internal $k$-uniform recycling and $k$-uniform product production in $\Gamma$ 
(again after the same bundlewise rescaling).

Because mass-action vector fields are polynomial and determined by the reaction 
monomials and their coefficients, the term-by-term agreement implies that the change 
of variables $[X_{B,k}]^N = s_B [X_{B,k}]$ conjugates the ODE of $\Gamma$ to that 
of $\Gamma^N$.
Since the vector fields are conjugate under a linear (diagonal, positive) change of 
coordinates, positive steady states correspond bijectively under this rescaling. 
A point $x^* \in \mathbb{R}^M_{>0}$ is a steady state of $\Gamma$ if and only if 
the rescaled point $(s_B x^*_{B,k})$ is a steady state of $\Gamma^N$.
Theorem~\ref{thm:BP-implementation} establishes that positive steady states of any 
CRN satisfying (W1)--(W6) and (R1) correspond to BP fixed points on the reconstructed 
factor graph. Applying this to both $\Gamma$ and $\Gamma^N$:
\begin{itemize}
\item Any positive steady state $x^*$ of $\Gamma$ corresponds to a BP fixed point on 
$\mathrm{FG}(\Gamma)$;
\item The rescaled steady state $(s_B x^*_{B,k})$ of $\Gamma^N$ corresponds to the 
same BP fixed point (since rescaling messages by positive constants does not change 
BP fixed-point status).
\end{itemize}

Thus steady states of $\Gamma$ determine BP fixed points on $\mathrm{FG}(\Gamma)$, 
with factor tables $\{\psi_J\}$ recovered from the reaction coefficients.
\end{proof}

\end{document}
